\section{Tính khách quan của khoa học}

\

Một mẩu thông tin đúng mà đa số mọi người công nhận với nhau là đúng được gọi là một \defText{khẳng định} hay \defText{sự thật}. Các khẳng định có thể là:
\begin{itemize}
    \item ``Mọi vật có sự sống cuối cùng đều sẽ chết.'';
    \item ``Hình vuông là tứ giác có bốn cạnh bằng nhau và bốn góc vuông.'';
    \item ``Nước đá là chất rắn, còn nước ở nhiệt độ thông thường là chất lỏng.''.
\end{itemize}

Trong trường hợp mà mẩu thông tin chưa được công nhận hay được kiểm chứng, thì sao? Nếu chúng ta có thể kiểm tra được tính chính xác của nó, thì chúng ta sẽ gọi nó là một \defText{giả thuyết}. Ví dụ:
\begin{itemize}
    \item ``Có sự sống ngoài Trái Đất.'';
    \item ``Loại thuốc $A$ này chữa được xoang.''.
\end{itemize}

Cần phải để ý rằng không có mệnh đề nào trong những ví dụ được đưa ra phụ thuộc vào quan điểm cá nhân của tác giả, của bạn đọc, hay của người phát biểu ra nó. Cho nên, những câu như:
\begin{itemize}
    \item ``Anh-xtanh\footnote{Albert Einstein (1879 -- 1955).} là một trong những nhà bác học vĩ đại nhất của thế kỉ XX.'', hay
    \item ``Tôi thích màu xanh của bầu trời hơn màu đỏ của hoàng hôn.''
\end{itemize}
đều không phải là khẳng định, mặc dù chúng có thể được coi là mệnh đề và xét tính đúng sai theo quan niệm của lô-gích diễn dịch.